\documentclass[addpoints, 11pt]{exam}
\usepackage[a4paper, left=20mm, right=20mm, top=26mm, bottom=30mm, headsep=6mm, footskip=12mm, headheight=15pt]{geometry}
\usepackage[T1]{fontenc}
\usepackage[utf8]{inputenc}
\usepackage{lmodern}
\usepackage{textcomp}
\usepackage[french]{babel}
\usepackage{xcolor}
\usepackage{listings}
\usepackage{graphicx}
\usepackage{array}
\usepackage{amsmath}
\usepackage{eurosym}
\usepackage{needspace}
\usepackage[most]{tcolorbox}
\usepackage{tikz}
\usetikzlibrary{positioning,arrows.meta}


% --- Code Python : même style que le sujet, en police à chasse fixe ---
\definecolor{codegreen}{rgb}{0,0.6,0}
\definecolor{codegray}{rgb}{0.5,0.5,0.5}
\definecolor{codepurple}{rgb}{0.58,0,0.82}
\definecolor{backcolour}{rgb}{0.95,0.95,0.92}
\lstdefinestyle{mystyle}{
    backgroundcolor=\color{backcolour},
    commentstyle=\color{codegreen},
    keywordstyle=\color{magenta},
    numberstyle=\tiny\color{codegray},
    stringstyle=\color{codepurple},
    basicstyle=\footnotesize\ttfamily,
    columns=fullflexible,
    breakatwhitespace=false,
    breaklines=true,
    postbreak=\mbox{\textcolor{codegray}{$\hookrightarrow$}\space},
    captionpos=b,
    keepspaces=true,
    numbers=left,
    numbersep=5pt,
    xleftmargin=1.5em,
    showspaces=false,
    showstringspaces=false,
    showtabs=false,
    tabsize=4,
    upquote=true
}
\lstset{style=mystyle, language=Python, extendedchars=true,
  literate={à}{{\`a}}1 {â}{{\^a}}1 {ä}{{\"a}}1 {ç}{{\c{c}}}1
           {é}{{\'e}}1 {è}{{\`e}}1 {ê}{{\^e}}1 {ë}{{\"e}}1
           {î}{{\^i}}1 {ï}{{\"i}}1 {ô}{{\^o}}1 {ö}{{\"o}}1
           {ù}{{\`u}}1 {û}{{\^u}}1 {ü}{{\"u}}1 {œ}{{\oe}}2
           {É}{{\'E}}1 {È}{{\`E}}1 {À}{{\`A}}1 {Ç}{{\c{C}}}1
           {°}{{\textdegree}}1 {€}{{\euro}}1 {…}{{\dots}}1
           {«}{{\guillemotleft}}1 {»}{{\guillemotright}}1 {·}{{\textperiodcentered}}1}
% Terminal (commandes Linux) : sans couleur de mots-clés
\lstdefinestyle{terminal}{style=mystyle, language={}, numbers=none, xleftmargin=0.5em}
% Sortie d'un programme (ce qui s'affiche à l'écran)
\lstdefinestyle{sortie}{style=mystyle, language={}, numbers=none, backgroundcolor=\color{white}, frame=single, rulecolor=\color{black!60}, xleftmargin=0.5em}

% --- Boîtes du corrigé (lisibles en noir et blanc) ---
\newtcolorbox{reponse}[1][Réponse attendue]{enhanced, breakable, colback=white, colframe=black, boxrule=1.2pt, arc=0pt,
  fonttitle=\bfseries, coltitle=white, colbacktitle=black, title=#1, left=2mm, right=2mm, top=1mm, bottom=1mm, before skip=6pt, after skip=6pt}
\newtcolorbox{erreur}[1][Erreur fréquente]{enhanced, breakable, colback=black!6, colframe=black!55, boxrule=0.6pt, arc=2pt,
  borderline west={3pt}{0pt}{black!70}, fonttitle=\bfseries, coltitle=black, colbacktitle=black!15, title=#1, left=2mm, right=2mm, top=1mm, bottom=1mm, before skip=6pt, after skip=6pt}
\newtcolorbox{astuce}[1][À retenir]{enhanced, breakable, colback=white, colframe=black!55, boxrule=0.6pt, arc=2pt,
  fonttitle=\bfseries, coltitle=black, colbacktitle=black!8, title=#1, left=2mm, right=2mm, top=1mm, bottom=1mm, before skip=6pt, after skip=6pt}
\newtcolorbox{remarque}[1][Remarque sur l'énoncé]{enhanced, breakable, colback=white, colframe=black!70, boxrule=0.6pt, arc=0pt,
  borderline={0.4pt}{1.2pt}{black!70, dashed}, fonttitle=\bfseries, coltitle=black, colbacktitle=white, title=#1, left=2mm, right=2mm, top=1mm, bottom=1mm, before skip=6pt, after skip=6pt}
\newcommand{\enonce}[1]{\textit{#1}\par\smallskip}
\newcommand{\methode}{\par\smallskip\noindent\textbf{Méthode.}\ }
\newcommand{\explication}{\par\smallskip\noindent\textbf{Explication.}\ }
\newcommand{\code}[1]{\texttt{#1}}
% Coupure de ligne possible après un tiret bas (noms de variables longs)
\DeclareRobustCommand{\_}{\textunderscore\allowbreak}
% Un code court ne se coupe pas en bas de page
\AddToHook{env/lstlisting/before}{\needspace{10\baselineskip}}
\setlength{\parindent}{0pt}
\setlength{\parskip}{3pt}
\renewcommand{\arraystretch}{1.25}

\title{Corrigé --- CC 1re NSI - Fonctions et listes (2025)}
\author{Babash - NSI Tani Malandi et Johanne Bakalara}
\date{Décembre 2025, durée : 1 h 30}

\begin{document}
\pagestyle{headandfoot}
\runningheader{}{}{\footnotesize Corrigé --- Fonctions et listes (2025)}
\firstpageheader{}{}{}
\runningfooter{}{\footnotesize Babash - NSI Tani Malandi et Johanne Bakalara \textperiodcentered{} CC BY-SA 4.0 \textperiodcentered{} page \thepage}{}
\firstpagefooter{}{\footnotesize Babash - NSI Tani Malandi et Johanne Bakalara \textperiodcentered{} CC BY-SA 4.0 \textperiodcentered{} page \thepage}{}

\maketitle
\thispagestyle{headandfoot}

\begin{astuce}[Comment lire ce corrigé]
Pour chaque question : la question est rappelée en italique, puis la \textbf{réponse attendue} est encadrée en noir, puis l'\textbf{explication} montre la méthode, étape par étape.
Les encadrés gris « Erreur fréquente » signalent les pièges les plus courants. Essaie d'abord de refaire la question seul·e, puis compare ta démarche avec celle du corrigé.
\end{astuce}

\section*{Barème}
\begin{center}
  \gradetable[h]
\end{center}

\begin{questions}

% =====================================================================
\section{Question de cours (fonctions)}

\question[1\half] \enonce{Soit la fonction \code{somme\_entre\_deux\_nombres}.}
%%% Bloc commun aux corrigés de fonctions (2024 et 2025) : questions sur somme_entre_deux_nombres
\begin{lstlisting}
def somme_entre_deux_nombres(a, b):
    somme = 0
    for i in range(a, b):
        somme += i
    return somme
\end{lstlisting}
\begin{parts}
\part \enonce{Quels sont les paramètres de la fonction ?}
\begin{reponse}
\code{a} et \code{b}.
\end{reponse}
Les \textbf{paramètres} sont les noms écrits entre parenthèses dans l'en-tête \code{def ...(a, b):}. Ce sont des variables qui reçoivent les valeurs données lors de l'appel.

\part \enonce{Quelles sont les variables locales de la fonction ?}
\begin{reponse}
\code{somme} et \code{i} (ainsi que les paramètres \code{a} et \code{b}, qui sont eux aussi locaux).
\end{reponse}
Une variable \textbf{locale} est créée à l'intérieur de la fonction et n'existe que pendant son exécution : après l'appel, \code{somme} et \code{i} ne sont plus accessibles depuis le programme principal. \code{somme} est créée ligne 2, \code{i} est créée par la boucle \code{for} ligne 3.

\part \enonce{Quelles sont les valeurs de retour de la fonction ?}
\begin{reponse}
Une seule valeur de retour : \code{somme}, c'est-à-dire la somme des entiers de \code{a} (inclus) à \code{b} (exclu).
\end{reponse}
La valeur de retour est ce qui suit le mot-clé \code{return}. Attention : \code{range(a, b)} s'arrête \textbf{avant} \code{b} ; \code{b} n'est pas ajouté.
\end{parts}


\question[1] \enonce{Écrire un appel de \code{somme\_entre\_deux\_nombres} pour \code{a=10} et \code{b=15}, en stockant les valeurs de retour dans des variables.}
\begin{reponse}
\begin{lstlisting}
resultat = somme_entre_deux_nombres(10, 15)
\end{lstlisting}
\code{resultat} vaut \textbf{60}.
\end{reponse}
\explication
\begin{enumerate}
  \item On écrit le nom de la fonction, puis les \textbf{arguments} entre parenthèses, dans l'ordre des paramètres : 10 va dans \code{a}, 15 va dans \code{b}.
  \item La fonction renvoie \textbf{une} valeur : on la range dans \textbf{une} variable avec \code{=}. Sans variable, la valeur renvoyée serait perdue.
\end{enumerate}
\textbf{Trace de l'appel :}
\begin{center}
\begin{tabular}{|c|c|c|}
\hline
Tour & \code{i} & \code{somme} \\ \hline
avant & & 0 \\
1 & 10 & 10 \\
2 & 11 & 21 \\
3 & 12 & 33 \\
4 & 13 & 46 \\
5 & 14 & 60 \\ \hline
\end{tabular}
\end{center}
\code{range(10, 15)} donne 10, 11, 12, 13, 14 : 15 est exclu.
\begin{erreur}
Écrire \code{def somme\_entre\_deux\_nombres(10, 15)} : le mot \code{def} sert à \textbf{définir}, pas à appeler. Autre erreur : appeler sans ranger le résultat : la valeur 60 est calculée puis perdue.
\end{erreur}

\question[4] \enonce{Associer à chaque numéro du schéma un terme : appeler, arguments, valeur de retour, corps, en-tête, définir, paramètre, nom.}
\begingroup
\tikzset{every picture/.append style={scale=0.78, transform shape}}
\begin{center}
\begin{tikzpicture}[font=\small,
  num/.style={draw,minimum width=0.7cm,minimum height=0.55cm,inner sep=1pt},
  lab/.style={anchor=west,inner sep=2pt},
  box/.style={draw=cyan!60!blue,rounded corners=6pt,very thick,inner sep=7pt}]
\node[box,anchor=center,align=left,font=\ttfamily\small] (g) at (0,0) {%
\textcolor{blue!70!black}{def} nbrCaractere (ch):\\
\hspace*{1.4em}i=0\\
\hspace*{1.4em}\textcolor{orange!80!black}{for} j \textcolor{orange!80!black}{in} ch:\\
\hspace*{2.8em}i+=1\\
\hspace*{1.4em}\textcolor{brown!80!black}{return} (i)};
\node[box,font=\ttfamily\small] (d) at (8,0) {nbrCaractere ('Alex')};
\node[num] (n1) at (-3,3.3) {1}; \node[lab] at (n1.east) {une fonction};
\node[num] (n2) at (5.6,1.6) {2}; \node[lab] at (n2.east) {une fonction};
\node[num] (n6) at (-1.6,2.3) {6}; \node[lab] at (n6.east) {de la fonction};
\node[num] (n8) at (2.6,2.3) {8};
\node[num] (n3) at (-6.6,1.0) {3}; \node[lab] (t3) at (n3.east) {de la fonction};
\node[num] (n4) at (-6.6,0) {4}; \node[lab] (t4) at (n4.east) {de la fonction};
\node[num] (n5) at (-6.6,-1.0) {5};
\node[num] (n7) at (8.3,-1.8) {7}; \node[lab] at (n7.east) {de la fonction};
\draw[->,thick] (t3.east) -- (-2.5,1.0);
\draw[->,thick] (t4.east) -- (-2.5,0);
\draw[->,thick] (n5.east) -- (-2.5,-1.0);
\draw[->,thick] (n6.south) -- (-0.6,1.15);
\draw[->,thick] (n8.south) -- (1.5,1.15);
\draw[->,thick] (n7.north) -- (8.3,-0.45);
\end{tikzpicture}
\end{center}

\endgroup
\begin{reponse}
\begin{tabular}{cl@{\qquad}cl}
1 & \textbf{définir} une fonction & 5 & \textbf{valeur de retour} \\
2 & \textbf{appeler} une fonction & 6 & \textbf{nom} de la fonction \\
3 & \textbf{en-tête} de la fonction & 7 & \textbf{arguments} de la fonction \\
4 & \textbf{corps} de la fonction & 8 & \textbf{paramètre} \\
\end{tabular}
\end{reponse}
\methode\ Commence par les mots que tu reconnais sûrement, puis termine par élimination.
\begin{enumerate}
  \item Le schéma a deux parties. À gauche, le bloc \code{def} est la \textbf{définition}.\\ À droite, la ligne \code{nbrCaractere('Alex')} est l'\textbf{appel}. Le 1 est au-dessus de la définition : « \textbf{définir} une fonction » ; le 2 est au-dessus de l'appel : « \textbf{appeler} une fonction ».
  \item Le 6 pointe vers le mot \code{nbrCaractere} : c'est le \textbf{nom}. Le 8 pointe vers \code{(ch)} : c'est le \textbf{paramètre}.
  \item Le 3 pointe vers la première ligne, celle du \code{def} : l'\textbf{en-tête}. Le 4 pointe vers les lignes décalées en dessous : le \textbf{corps}.
  \item Le 5 pointe vers la ligne \code{return (i)} : ce qui est renvoyé est la \textbf{valeur de retour}.
  \item Le 7 pointe vers \code{'Alex'}, la valeur donnée lors de l'appel : c'est l'\textbf{argument}.
\end{enumerate}
\begin{erreur}
Confondre \textbf{paramètre} et \textbf{argument}. Le paramètre (\code{ch}) est le nom écrit dans la définition ; l'argument (\code{'Alex'}) est la valeur concrète donnée à l'appel. Lors de l'appel, \code{ch} reçoit \code{'Alex'}, et la fonction renvoie 4 (le nombre de caractères).
\end{erreur}


\question[3] \enonce{Dans le code du jeu « devine le nombre », entourer les définitions des fonctions (en entier) et souligner les appels.}
\begin{reponse}
\textbf{Définitions à entourer} (chacune en entier, de la ligne \code{def} à la dernière ligne indentée) :
\begin{itemize}
  \item \code{demander\_nombre} : lignes 3 à 5 ;
  \item \code{aide} : lignes 7 à 11.
\end{itemize}
\textbf{Appels à souligner} :
\begin{itemize}
  \item ligne 20 : \code{demander\_nombre()} ;
  \item ligne 21 : \code{aide(rep, nbHasard)}.
\end{itemize}
On accepte aussi de souligner les appels de fonctions prédéfinies : \code{int(...)} et \code{input(...)} (ligne 4), \code{print(...)} (lignes 9, 11, 24, 26), \code{random.randint(1, 100)} (ligne 15).
\end{reponse}
\textbf{Le code annoté :}
\begin{lstlisting}[escapechar=@, numbers=none]
import random

@\fbox{\begin{minipage}{0.75\linewidth}\ttfamily\footnotesize\textcolor{magenta}{def} demander\_nombre():\\
\hspace*{2em}rep = int(input("quelle est votre choix ?: "))\\
\hspace*{2em}\textcolor{magenta}{return} rep\end{minipage}}@

@\fbox{\begin{minipage}{0.75\linewidth}\ttfamily\footnotesize\textcolor{magenta}{def} aide(rep, nbHasard):\\
\hspace*{2em}\textcolor{magenta}{if} rep<nbHasard:\\
\hspace*{4em}print("+")\\
\hspace*{2em}\textcolor{magenta}{if} rep>nbHasard:\\
\hspace*{4em}print("-")\end{minipage}}@

essai = 11
#on crée une variable qui va stocker le nombre à trouver .
nbHasard = random.randint(1 , 100)
rep = 0

while rep != nbHasard and essai > 0:
    essai = essai - 1
    rep = @\underline{demander\_nombre()}@
    @\underline{aide(rep, nbHasard)}@

if essai > 0 :
    print("Bravo, vous avez trouvé la bonne réponse !")
else :
    print("Vous avez perdu !")
\end{lstlisting}

\methode
\begin{enumerate}
  \item Une \textbf{définition} commence par le mot-clé \code{def}, suivi du nom, des parenthèses et des deux-points. Elle comprend \textbf{tout le corps} : toutes les lignes indentées sous le \code{def}. On l'entoure en entier.
  \item Un \textbf{appel} est le nom d'une fonction suivi de parenthèses, \textbf{sans} \code{def} devant : c'est le moment où la fonction s'exécute.
  \item Une fonction définie mais jamais appelée ne fait rien. Ici, les deux fonctions sont appelées à chaque tour de la boucle \code{while}.
\end{enumerate}
\begin{erreur}
Souligner la ligne \code{def demander\_nombre():} comme un appel. La ligne \code{def} \textbf{définit} la fonction (elle explique ce qu'elle fera) ; elle ne l'exécute pas.
\end{erreur}


% =====================================================================
\section{Listes Python}

\question[1\half] \enonce{Manipulation de base.}
\begin{parts}
\part \enonce{Comment déclare-t-on une liste vide en Python ?}
\begin{reponse}
\lstinline!ma_liste = []!
\end{reponse}
Deux crochets sans rien entre eux. Variante : \code{ma\_liste = list()}.

\part \enonce{Quel est le rôle de la fonction \code{len()} appliquée à une liste ?}
\begin{reponse}
\code{len()} renvoie la \textbf{longueur} de la liste, c'est-à-dire son nombre d'éléments.
\end{reponse}
Exemple : \code{len(["Alice", "Bob", "Charlie"])} vaut 3 ; \code{len([])} vaut 0. Le dernier élément a donc l'indice \code{len(liste) - 1}.

\part \enonce{Quelle méthode ajoute un élément à la fin d'une liste existante ?}
\begin{reponse}
La méthode \code{append} : \lstinline!ma_liste.append(element)!
\end{reponse}
Une \textbf{méthode} est une fonction attachée à un objet : on l'écrit après un point. Exemple : si \code{L = []}, après \code{L.append(3)}, \code{L} vaut \code{[3]}.
\end{parts}

\question[2] \enonce{Écrire deux boucles qui parcourent \code{noms = ["Alice", "Bob", "Charlie"]}.}
\begin{parts}
\part \enonce{Un parcours par valeur.}
\begin{reponse}
\begin{lstlisting}
for nom in noms:
    print(nom)
\end{lstlisting}
\end{reponse}
La variable \code{nom} prend directement chaque \textbf{valeur} de la liste : \code{"Alice"}, puis \code{"Bob"}, puis \code{"Charlie"}.

\part \enonce{Un parcours par index.}
\begin{reponse}
\begin{lstlisting}
for i in range(len(noms)):
    print(noms[i])
\end{lstlisting}
\end{reponse}
La variable \code{i} prend chaque \textbf{indice} (la position) : \code{len(noms)} vaut 3, donc \code{range(3)} donne 0, 1, 2. On lit l'élément avec \code{noms[i]}.
\begin{center}
\begin{tabular}{|c|c|c|}
\hline
Tour & \code{i} & \code{noms[i]} \\ \hline
1 & 0 & \code{"Alice"} \\
2 & 1 & \code{"Bob"} \\
3 & 2 & \code{"Charlie"} \\ \hline
\end{tabular}
\end{center}
Les deux boucles affichent la même chose. Le parcours par index est utile quand on a besoin de la position (par exemple pour comparer un élément avec son voisin).
\begin{erreur}
Écrire \code{for i in noms: print(noms[i])} : \code{i} contient \code{"Alice"}, pas un nombre, et \code{noms["Alice"]} provoque une \code{TypeError} (« les indices de liste doivent être des entiers »).
\end{erreur}
\end{parts}

% =====================================================================
\section{Analyse de code}

\question[7] \enonce{Prédire la sortie de chacun des codes suivants.}
\begin{parts}
\part
\begin{lstlisting}
def operation_simple(n):
    return n + 3

def processus_A(x):
    y = operation_simple(x)
    return y * 4

valeur1 = 7
valeur2 = 2
resultat1 = processus_A(valeur1)
resultat2 = processus_A(valeur2)
print(resultat1)
print(resultat2)
\end{lstlisting}
\begin{reponse}
\begin{lstlisting}[style=sortie]
40
20
\end{lstlisting}
\end{reponse}
\code{processus\_A(7)} appelle \code{operation\_simple(7)}, qui renvoie $7 + 3 = 10$ ; puis $10 \times 4 = 40$. De même, \code{processus\_A(2)} : $(2 + 3) \times 4 = 20$. Une fonction peut en appeler une autre : on calcule d'abord l'appel le plus intérieur.

\part
\begin{lstlisting}
def traitement_B(L):
    element_cle = L[0]
    for e in L:
        if e < element_cle:
            element_cle = e
    return element_cle

nombres = [1, 50, 206, 9, 17]
print(traitement_B(nombres))
\end{lstlisting}
\begin{reponse}
\begin{lstlisting}[style=sortie]
1
\end{lstlisting}
La fonction renvoie le \textbf{minimum} de la liste.
\end{reponse}
\begin{center}
\begin{tabular}{|c|c|c|c|}
\hline
Tour & \code{e} & \code{e < element\_cle} ? & \code{element\_cle} \\ \hline
avant & & & 1 \\
1 & 1 & $1 < 1$ : faux & 1 \\
2 & 50 & faux & 1 \\
3 & 206 & faux & 1 \\
4 & 9 & faux & 1 \\
5 & 17 & faux & 1 \\ \hline
\end{tabular}
\end{center}
\code{element\_cle} démarre avec le premier élément, puis est remplacé chaque fois qu'on trouve plus petit. Ici le premier élément est déjà le plus petit.

\part
\begin{lstlisting}
def fonction_C(liste):
    resultat = 0
    for val in liste:
        resultat = resultat + val
    return resultat

donnees = [1, 2, 3, 4, 5]
print(fonction_C(donnees))
\end{lstlisting}
\begin{reponse}
\begin{lstlisting}[style=sortie]
15
\end{lstlisting}
La fonction renvoie la \textbf{somme} des éléments.
\end{reponse}
\code{resultat} prend successivement les valeurs 0, 1, 3, 6, 10, 15 (on ajoute 1, puis 2, puis 3, puis 4, puis 5).

\part
\begin{lstlisting}
def transformation_D(L):
    L_nouvelle = []
    for i in range(len(L) - 1):
        paire = L[i] + L[i+1]
        L_nouvelle.append(paire)
    return L_nouvelle

original = [10, 3, 2, 5, 2, 8]
resultat = transformation_D(original)
print(resultat)
\end{lstlisting}
\begin{reponse}
\begin{lstlisting}[style=sortie]
[13, 5, 7, 7, 10]
\end{lstlisting}
Chaque élément est la somme de deux voisins.
\end{reponse}
\begin{center}
\begin{tabular}{|c|c|c|c|l|}
\hline
\code{i} & \code{L[i]} & \code{L[i+1]} & \code{paire} & \code{L\_nouvelle} \\ \hline
0 & 10 & 3 & 13 & \code{[13]} \\
1 & 3 & 2 & 5 & \code{[13, 5]} \\
2 & 2 & 5 & 7 & \code{[13, 5, 7]} \\
3 & 5 & 2 & 7 & \code{[13, 5, 7, 7]} \\
4 & 2 & 8 & 10 & \code{[13, 5, 7, 7, 10]} \\ \hline
\end{tabular}
\end{center}
La liste a 6 éléments : \code{range(len(L) - 1)} donne 0 à 4. On s'arrête à 4 car \code{L[i+1]} doit exister : avec \code{i = 5}, \code{L[6]} provoquerait une \code{IndexError}. Le résultat a donc un élément de moins que la liste de départ.

\part
\begin{lstlisting}
def analyse_E(L, cible):
    for element in L:
        if element == cible:
            return True
    return False

liste_a = ["a", "b", "c"]
print(analyse_E(liste_a, "d"))
\end{lstlisting}
\begin{reponse}
\begin{lstlisting}[style=sortie]
False
\end{lstlisting}
La fonction teste si \code{cible} est dans la liste.
\end{reponse}
Aucun élément n'est égal à \code{"d"} : le \code{return True} n'est jamais atteint. La boucle se termine, puis on exécute \code{return False}, qui est \textbf{après} la boucle (même indentation que le \code{for}).
\begin{erreur}
Croire que la fonction renvoie \code{False} dès le premier élément différent. Le \code{return False} est en dehors de la boucle : on ne l'atteint qu'après avoir regardé \textbf{tous} les éléments.
\end{erreur}

\part
\begin{lstlisting}
def methode_F(L, v):
    compteur = 0
    for element in L:
        if element == v:
            compteur = compteur + 1
    return compteur

lettres = ["a", "b", "a", "c", "a", "b"]
print(methode_F(lettres, "b"))
\end{lstlisting}
\begin{reponse}
\begin{lstlisting}[style=sortie]
2
\end{lstlisting}
La fonction compte le nombre d'apparitions de \code{v} dans la liste.
\end{reponse}
\code{"b"} est aux indices 1 et 5 : le compteur passe de 0 à 1, puis à 2.

\part
\begin{lstlisting}
nombres = [10, 100, 1000]
chaine = ""
for i in range(len(nombres)):
    chaine = chaine + str(nombres[i]) + "-"
print(chaine)
\end{lstlisting}
\begin{reponse}
\begin{lstlisting}[style=sortie]
10-100-1000-
\end{lstlisting}
\end{reponse}
\begin{center}
\begin{tabular}{|c|c|l|}
\hline
\code{i} & \code{str(nombres[i])} & \code{chaine} \\ \hline
avant & & \code{""} \\
0 & \code{"10"} & \code{"10-"} \\
1 & \code{"100"} & \code{"10-100-"} \\
2 & \code{"1000"} & \code{"10-100-1000-"} \\ \hline
\end{tabular}
\end{center}
\code{str()} transforme le nombre en chaîne ; le \code{+} entre deux chaînes les \textbf{colle} (on parle de \textbf{concaténation}). Un tiret est ajouté après chaque nombre, y compris le dernier. Le \code{print} est en dehors de la boucle : il n'affiche qu'une fois, à la fin.
\begin{erreur}
Oublier \code{str()} : \code{"" + 10} provoque une \code{TypeError}, car Python n'additionne pas une chaîne et un entier.
\end{erreur}
\end{parts}

\end{questions}

\end{document}
